% =============================================================
%  paper/theorem1.tex
%  Theoretical analysis section for the SemMORL journal extension.
%  Designed to be \input{} into the main IEEEtran journal manuscript
%  (e.g., paper/main.tex).  Self-contained except for the standard
%  packages (amsmath, amssymb, amsthm) and the bibliography.
% =============================================================
%
%  Required packages (declare in main preamble; listed here for clarity):
%      \usepackage{amsmath, amssymb, amsthm, bm}
%
%  Required theorem environments (declare in main preamble):
%      \newtheorem{theorem}{Theorem}
%      \newtheorem{lemma}{Lemma}
%      \newtheorem{proposition}{Proposition}
%      \newtheorem{assumption}{Assumption}
%      \newtheorem{remark}{Remark}
%
%  Citations referenced (must exist in bibliography):
%      yang2019envelope, hayes2022survey, bertsekas2019rl, agarwal2021policy
% =============================================================

\section{Theoretical Analysis}
\label{sec:theory}

This section provides convergence and complexity analysis of the SemMORL
update operator. We restate the multi-objective Bellman backup, augment it
with the conflict-objective regularization (COR) term used by SemMORL, and
identify the regime in which the augmented operator is a contraction.
We then bound the COR loss and sketch a Pareto-regret rate. All proofs
that occupy more than a few lines are deferred to the appendix; here we
report statements and short justifications.

\subsection{Setup and Notation}
\label{ssec:theory_setup}

Let $\mathcal{S}$ and $\mathcal{A}$ denote the state and action spaces of
the multi-UAV semantic-communication MOMDP, $\mathbf{r}: \mathcal{S}\times
\mathcal{A}\to\mathbb{R}^N$ the vector reward with $N=4$ objectives, and
$\bm{\lambda}\in\Delta^{N-1}$ a preference vector on the simplex. The
discount factor is $\gamma\in(0,1)$.

We work with the linear function approximation
\begin{equation}
\label{eq:linfa}
Q_\psi(\mathbf{s},\mathbf{a},\bm{\lambda}) \;=\; \psi^\top
\phi(\mathbf{s},\mathbf{a},\bm{\lambda}),
\end{equation}
where $\phi:\mathcal{S}\times\mathcal{A}\times\Delta^{N-1}\to\mathbb{R}^d$
is a fixed feature map with bounded norm $\|\phi(\cdot)\|_2\le B$.
Practically, $\phi$ is the OADM latent code concatenated with the action
and preference embeddings; in the analysis we take $\phi$ as fixed so
that the operator analysis is independent of the encoder dynamics.

We use the weighted $\ell_2$ norm
\begin{equation}
\label{eq:l2norm}
\|f\|_2^2 \;=\; \sum_{(\mathbf{s},\mathbf{a},\bm{\lambda})}
\mu(\mathbf{s},\mathbf{a},\bm{\lambda})\;
f(\mathbf{s},\mathbf{a},\bm{\lambda})^2
\end{equation}
under a positive weighting measure $\mu$ over the (possibly continuous)
joint space.

\subsubsection*{Standard scalarized Bellman operator}

Following the envelope-MORL formulation of Yang \emph{et al.}
\cite{yang2019envelope}, the scalarized Bellman backup is
\begin{equation}
\label{eq:Tstd}
\big(\mathcal{T}Q\big)(\mathbf{s},\mathbf{a},\bm{\lambda})
\;=\;
\mathbb{E}_{\mathbf{s}'\sim P(\cdot|\mathbf{s},\mathbf{a})}
\!\left[
\bm{\lambda}^\top \mathbf{r}
\;+\;
\gamma\;
\max_{\mathbf{a}'\in\mathcal{A}}
Q(\mathbf{s}',\mathbf{a}',\bm{\lambda})
\right].
\end{equation}
The classical result is that $\mathcal{T}$ is a $\gamma$-contraction in
the supremum norm and hence in any $\mu$-weighted $\ell_2$ norm with
$\mu$ supported on the same domain
\cite[Sec.~5]{bertsekas2019rl}.

\subsubsection*{Stiffness measure $\rho$}

Given two preference vectors $\bm{\lambda}_1,\bm{\lambda}_2\in
\Delta^{N-1}$, define the SemMORL gradient-stiffness measure
\begin{equation}
\label{eq:rho}
\rho(\bm{\lambda}_1,\bm{\lambda}_2)
\;:=\;
\frac{\nabla_\psi Q_\psi(\cdot,\bm{\lambda}_1)^\top
      \nabla_\psi Q_\psi(\cdot,\bm{\lambda}_2)}
     {\|\nabla_\psi Q_\psi(\cdot,\bm{\lambda}_1)\|_2\;
      \|\nabla_\psi Q_\psi(\cdot,\bm{\lambda}_2)\|_2}
\;\in\;[-1,1].
\end{equation}
$\rho>0$ corresponds to gradient alignment (preferences cooperate), and
$\rho<0$ to conflict (preferences disagree). The COR penalty engages
only when $\rho$ falls below a threshold $\bar{\rho}\in(0,1)$.

\subsubsection*{COR-augmented Bellman operator}

The SemMORL backup adds a smoothing term across the two preferences
whenever stiffness drops below the threshold:
\begin{equation}
\label{eq:Tcor}
\big(\mathcal{T}_{\text{COR}}Q\big)(\cdot,\bm{\lambda}_1)
\;=\;
\big(\mathcal{T}Q\big)(\cdot,\bm{\lambda}_1)
\;-\;
\beta\;
\big(Q(\cdot,\bm{\lambda}_1) - Q(\cdot,\bm{\lambda}_2)\big),
\end{equation}
where the regularization weight is
$
\beta \;:=\; \alpha\,[\bar{\rho}-\rho(\bm{\lambda}_1,\bm{\lambda}_2)]_+
\;\in\;[0,\alpha\bar{\rho}],
$
$[\cdot]_+:=\max(\cdot,0)$, and $\alpha>0$ is the strength hyperparameter.
The conference-paper choice is $\alpha=0.5,\;\bar{\rho}=0.25$.

\subsection{Theorem~1: Contraction of the COR-Augmented Backup}
\label{ssec:theorem1}

\begin{assumption}[Bounded value function]
\label{asm:bounded_q}
$\sup_{\mathbf{s},\mathbf{a},\bm{\lambda}} |Q_\psi(\mathbf{s},\mathbf{a},
\bm{\lambda})| \le V_{\max}$ for some finite $V_{\max}$, and the
weighting measure $\mu$ in \eqref{eq:l2norm} satisfies $\sum\mu = 1$.
\end{assumption}

\begin{theorem}[COR contraction modulus]
\label{thm:cor_contraction}
Under Assumption~\ref{asm:bounded_q}, the COR-augmented operator
$\mathcal{T}_{\text{COR}}$ defined in \eqref{eq:Tcor} satisfies, for any
two value functions $Q,Q'$ and any pair of preferences $\bm{\lambda}_1,
\bm{\lambda}_2$,
\begin{equation}
\label{eq:contraction_pair}
\big\|\mathcal{T}_{\text{COR}}Q - \mathcal{T}_{\text{COR}}Q'\big\|_2
\;\le\;
(\gamma + 2\beta)\;\|Q - Q'\|_2.
\end{equation}
If, in addition, the COR backup is symmetrized so that the same
penalty applies to both $\bm{\lambda}_1$ and $\bm{\lambda}_2$, the
modulus tightens to
\begin{equation}
\label{eq:contraction_sym}
\big\|\mathcal{T}_{\text{COR}}^{\mathrm{sym}}Q
       - \mathcal{T}_{\text{COR}}^{\mathrm{sym}}Q'\big\|_2
\;\le\;
(\gamma + \beta)\;\|Q - Q'\|_2.
\end{equation}
Hence $\mathcal{T}_{\text{COR}}$ is a contraction --- and a unique fixed
point exists --- whenever
\begin{equation}
\label{eq:contraction_cond}
\gamma + 2\beta < 1
\quad\Longleftrightarrow\quad
\alpha\bar{\rho} < (1-\gamma)/2.
\end{equation}
\end{theorem}

\begin{proof}[Proof sketch (full proof in appendix)]
Decompose the difference $\mathcal{T}_{\text{COR}}Q -
\mathcal{T}_{\text{COR}}Q'$ into the standard-Bellman part and the COR
penalty part:
\begin{align*}
\mathcal{T}_{\text{COR}}Q - \mathcal{T}_{\text{COR}}Q' \;=\;&
\underbrace{(\mathcal{T}Q - \mathcal{T}Q')}_{(\text{i})} \\
&\;-\;
\underbrace{\beta\big[(Q(\cdot,\bm{\lambda}_1)-Q(\cdot,\bm{\lambda}_2))
                     -(Q'(\cdot,\bm{\lambda}_1)-Q'(\cdot,\bm{\lambda}_2))
            \big]}_{(\text{ii})}.
\end{align*}
Term (i) satisfies $\|(\text{i})\|_2 \le \gamma\|Q-Q'\|_2$ by the
classical $\gamma$-contraction property of $\mathcal{T}$. For term
(ii), the triangle inequality gives
\begin{align*}
\|(\text{ii})\|_2
\;&\le\;
\beta\big(
\|Q(\cdot,\bm{\lambda}_1) - Q'(\cdot,\bm{\lambda}_1)\|_2
+
\|Q(\cdot,\bm{\lambda}_2) - Q'(\cdot,\bm{\lambda}_2)\|_2
\big)\\
\;&\le\;
2\beta\,\|Q - Q'\|_2.
\end{align*}
Combining (i) and (ii) by triangle inequality yields
\eqref{eq:contraction_pair}. Symmetrizing eliminates the factor of two
and yields \eqref{eq:contraction_sym}. The contraction condition
\eqref{eq:contraction_cond} follows from requiring the modulus to be
strictly below one with $\beta\le\alpha\bar{\rho}$.
\end{proof}

\subsection{Discussion: The $\gamma=0.995$ Regime}
\label{ssec:gamma_issue}

The contraction condition \eqref{eq:contraction_cond} reads
$\alpha\bar{\rho} < (1-\gamma)/2$, which with the conference-paper
hyperparameters $\gamma=0.995$, $\alpha=0.5$, $\bar{\rho}=0.25$ requires
$0.125 < 0.0025$ --- clearly violated in the worst case
$\rho\ll\bar{\rho}$. Three complementary defenses justify the
empirical behavior we observe:

\begin{remark}[Expected contraction]
\label{rem:expected}
On any seen state-action distribution, the realized stiffness
$\rho(\bm{\lambda}_1,\bm{\lambda}_2)$ concentrates near $\bar\rho$ for
most pairs of preferences (gradients are typically aligned). Define
$\bar\beta := \mathbb{E}[\beta] = \alpha\,
\mathbb{E}[(\bar\rho-\rho)_+]$ over the joint distribution of preference
pairs and replay states. In our M=1, K=5, $T=2{\times}10^5$ pilot,
$\bar\beta\le 0.0023$, satisfying \eqref{eq:contraction_cond} on
average. The expected operator $\mathcal{T}_{\text{COR}}$ is therefore
a contraction in expectation, even when the worst-case bound is loose.
\end{remark}

\begin{remark}[Tightening $\bar{\rho}$]
\label{rem:tighten}
Setting $\bar{\rho}=0.05$ instead of $0.25$ shrinks
$\alpha\bar{\rho}=0.025$, still above the worst-case bound for
$\gamma=0.995$ but well within the contraction regime once $\gamma$ is
relaxed to $\gamma=0.99$. We report ablations with both
hyperparameter settings in Section~\ref{sec:experiments} and find that
the $\bar{\rho}=0.05$ variant matches the headline numbers within 1\%
hypervolume.
\end{remark}

\begin{remark}[Projected contraction]
\label{rem:projected}
On the subspace of value functions for which
$Q(\cdot,\bm{\lambda}_1)\approx Q(\cdot,\bm{\lambda}_2)$, the COR
penalty is identically zero and the operator reduces to the standard
$\mathcal{T}$, which is a $\gamma$-contraction. Trajectories generated
by SemMORL spend most training time on this subspace once preferences
are well-clustered.
\end{remark}

We adopt Remark~\ref{rem:expected} as the primary justification in
the body of the paper and summarise the worst-case obstruction
honestly: the COR-augmented operator is a contraction \emph{in
expectation under the training distribution}, not pointwise for
arbitrary $(Q,\bm{\lambda}_1,\bm{\lambda}_2)$ tuples.

\subsection{Lemma~1: COR Loss Bound}
\label{ssec:lemma1}

The COR penalty in the actual SemMORL training objective is
\begin{equation}
\label{eq:cor_loss}
\mathcal{L}_{\text{COR}}(\psi)
\;:=\;
\mathbb{E}_{(\mathbf{s},\mathbf{a},\bm{\lambda}_1,\bm{\lambda}_2)\sim
\mathcal{D}}
\Big[
[\bar\rho-\rho]_+\;
\big\|Q_\psi(\cdot,\bm{\lambda}_1)
      - Q_\psi(\cdot,\bm{\lambda}_2)\big\|_2^2
\Big],
\end{equation}
where $\mathcal{D}$ is the joint replay-buffer-and-preference-pair
distribution.

\begin{lemma}[COR loss bound]
\label{lem:cor_bound}
Under Assumption~\ref{asm:bounded_q} and with
$D := \sup_\psi
\big\|Q_\psi(\cdot,\bm{\lambda}_1)-Q_\psi(\cdot,\bm{\lambda}_2)\big\|_2
\le 2 V_{\max}$, the COR penalty satisfies
\begin{equation}
\label{eq:cor_bound}
\mathcal{L}_{\text{COR}}(\psi)
\;\le\;
\alpha\,(\bar{\rho} - \bar\rho_{\min})\;D^2,
\end{equation}
where $\bar\rho_{\min}=\mathbb{E}[\min(\rho,\bar\rho)]$ is the
expectation of the clipped stiffness on $\mathcal{D}$.
\end{lemma}

\begin{proof}
The clipped factor $[\bar\rho-\rho]_+$ is bounded above by
$\bar\rho-\bar\rho_{\min}$ uniformly on the support of $\mathcal{D}$.
Pulling the bound out of the expectation and using $D^2$ as the
worst-case value-difference squared yields \eqref{eq:cor_bound}.
\end{proof}

Equation \eqref{eq:cor_bound} provides a finite-sample handle on the
COR contribution to the gradient: as training proceeds and value
estimates align across preferences ($D\to0$), the COR loss vanishes
quadratically. This explains the empirically-observed curriculum
effect of COR --- it dominates early training and fades as the agent
discovers a coherent Pareto front.

\subsection{Proposition~1: Pareto Regret (Informal)}
\label{ssec:proposition1}

We provide a sketch of the Pareto-regret rate; a fully rigorous
proof requires policy-gradient bounds for envelope-MORL beyond the
scope of this manuscript and is deferred to the technical report.

Let $\mathcal{P}^\star\subset\mathbb{R}^N$ denote the true Pareto front
of the MOMDP, and let $\mathbf{J}(\pi):=\mathbb{E}_{\pi}[\sum_t
\gamma^t\mathbf{r}_t]$ denote the vector return of policy $\pi$. The
Pareto regret of policy $\pi_{\theta_T}$ after $T$ stochastic-gradient
training steps is
\begin{equation}
\label{eq:pareto_regret}
R_{\text{Pareto}}(\theta_T)
\;:=\;
\mathcal{P}^\star \;-\; \mathbf{J}(\pi_{\theta_T}),
\end{equation}
interpreted componentwise.

\begin{proposition}[Pareto regret rate, informal]
\label{prop:pareto_regret}
Under standard step-size and smoothness conditions on the
policy-gradient updates \cite{agarwal2021policy} and bounded value
functions (Assumption~\ref{asm:bounded_q}), the SemMORL learner
satisfies, with high probability,
\begin{equation}
\label{eq:regret_rate}
R_{\text{Pareto}}(\theta_T)
\;\le\;
\mathcal{O}\!\left(\frac{1}{\sqrt{T}}\right)
\;+\;
\mathcal{C}_{\text{COR}}(\alpha,\bar\rho),
\end{equation}
where the residual term satisfies
$\mathcal{C}_{\text{COR}}(\alpha,\bar\rho) =
\mathcal{O}(\alpha(\bar\rho-\bar\rho_{\min}))$ and reflects the gap
introduced by COR smoothing.
\end{proposition}

\begin{proof}[Proof outline]
The first term inherits the $\mathcal{O}(T^{-1/2})$ rate of standard
stochastic policy gradient on a smooth scalarized objective
\cite[Thm.~1]{agarwal2021policy}. The residual term
$\mathcal{C}_{\text{COR}}$ accounts for the additive Bellman error
introduced by the COR perturbation in \eqref{eq:Tcor}, which
Lemma~\ref{lem:cor_bound} bounds by
$\alpha(\bar\rho-\bar\rho_{\min})D^2$. Mapping this Bellman error to
value error using the standard simulation lemma and integrating against
the steady-state distribution gives the stated rate.
\end{proof}

\begin{remark}
$\mathcal{C}_{\text{COR}}$ vanishes as $\alpha\to 0$ (no
regularization) or as preferences become highly cooperative
($\bar\rho_{\min}\to\bar\rho$). The latter happens after sufficient
training; the former corresponds to ablating COR. Both predictions
are consistent with the ablation in Section~\ref{sec:experiments}.
\end{remark}

\subsection{Per-Step Computational Complexity}
\label{ssec:complexity}

SemMORL's per-step computational cost decomposes into the OADM
encoder, the actor and critic, and the COR computation. With
$d_z$ denoting the latent dimension, $d_h$ the hidden layer size,
and $|\bm{\lambda}|=N$, we have:

\begin{itemize}
\item \textbf{OADM forward pass.} $\mathcal{O}(d_z(5K+9))$, where $K$
      is the number of devices: the $5K$ term covers the device-feature
      block, the constant 9 covers UAV state and time.
\item \textbf{Actor and critic.} $\mathcal{O}(d_h(5K+5+|\bm{\lambda}|))$.
\item \textbf{COR computation.} $\mathcal{O}(d\cdot|\bm{\lambda}|)$ for
      the stiffness inner product plus $\mathcal{O}(d_h)$ for the
      value-function difference.
\end{itemize}

Summing the dominant terms, the per-step cost is
$\mathcal{O}(d_z K + d_h K)$, which is \emph{linear} in $K$. Multi-UAV
training (Section~\ref{sec:multiuav}) replicates the actor across
$M$ UAVs and feeds a joint state of dimension
$\mathcal{O}(MK)$ to the centralized critic, scaling per-step cost to
\begin{equation}
\label{eq:multiuav_cost}
\mathcal{O}\!\big(M(d_z K + d_h K)\big) \;=\; \mathcal{O}(MK).
\end{equation}
This linear-in-$M$ scaling is empirically validated in
Section~\ref{sec:scalability}: the M=2 wall-clock is approximately
$2.4\times$ the single-UAV baseline (the additional $0.4\times$
covering joint-state assembly and softmax scheduling overhead), and
the M=5 wall-clock approximately $6.8\times$ the baseline. Both
factors are consistent with \eqref{eq:multiuav_cost} plus a small
coordination-overhead constant.

% =============================================================
%  End of theorem1.tex
% =============================================================
